\section{Tujuan Buku Ajar}

Buku ajar ini dirancang sebagai panduan komprehensif untuk menguasai \logic{Logika Matematika}, sejalan dengan filosofi pendidikan berbasis outcome yang menekankan pencapaian kompetensi terukur \cite{spady1994,openlogic,rosen2019,ramdhania2023}. Sebagaimana dikembangkan oleh sumber-sumber terkemuka seperti Open Logic Project dan Rosen, materi logika formal disusun secara sistematis dari dasar hingga aplikasi praktis dalam ilmu komputer. Buku ini bertujuan memfasilitasi mahasiswa Teknik Informatika agar mencapai learning outcomes yang telah ditetapkan dalam kurikulum Program Studi.

Pentingnya logika matematika bagi mahasiswa informatika tidak dapat diabaikan, mengingat keterampilan berpikir logis menjadi fondasi bagi pemrograman, desain algoritma, kecerdasan buatan, dan verifikasi perangkat lunak \cite{levin_discrete,stanford_cs103}. Mata kuliah CS103 Stanford dan Discrete Mathematics Levin menempatkan logika proposisional dan predikat sebagai prasyarat untuk memahami struktur diskret yang mendasari komputasi. Penguasaan logika memungkinkan mahasiswa menganalisis kebenaran program, merancang kondisi dalam kode, dan memahami dasar-dasar reasoning otomatis dalam sistem cerdas.

Buku ini mengikuti alur standar textbook logika formal yang diadopsi secara internasional \cite{forallx,copi2014}. Seperti \textit{forall x: Calgary} dan \textit{Introduction to Logic} Copi, urutan materi meliputi proposisi, operator logika, tabel kebenaran, ekuivalensi, logika predikat dengan kuantifier, serta metode pembuktian. Alur ini memastikan mahasiswa membangun pemahaman bertahap sebelum beralih ke topik yang lebih kompleks seperti induksi matematika dan logic programming.

Tujuan spesifik buku ini adalah:

\begin{enumerate}
  \item Memberikan pemahaman mendalam tentang proposisi dan operator logika
  \item Mengembangkan kemampuan menerjemahkan pernyataan ke bentuk logika formal
  \item Membangun keterampilan membangun bukti formal menggunakan aturan inferensi
  \item Memperkenalkan logika predikat dan induksi matematika
  \item Mengenalkan dasar-dasar Prolog sebagai penerapan logic programming
  \item Memfasilitasi pencapaian CPL dan CPMK yang telah ditetapkan
\end{enumerate}

Setelah mempelajari buku ini secara menyeluruh, mahasiswa diharapkan mampu mencapai kompetensi berikut \cite{learn_prolog}. Pengenalan Prolog sebagai bahasa logic programming memperlihatkan penerapan praktis konsep logika dalam penulisan fakta, aturan, dan query. Kemampuan-kemampuan tersebut berkontribusi langsung terhadap CPL Program Studi Teknik Informatika dalam aspek pengetahuan dan keterampilan analitis.

\begin{itemize}
  \item Membuat tabel kebenaran dan menganalisis tautologi
  \item Menerjemahkan pernyataan bahasa alami ke simbol logika
  \item Membangun bukti menggunakan modus ponens, modus tollens, dan induksi
  \item Menulis fakta dan aturan sederhana dalam Prolog
\end{itemize}
